\section{Problem Statement and Network Infrastructure Architecture}

\subsection{Operational Network Challenges}
The deployment and maintenance of Analysis of TCP and UDP Protocol Performance in complex organizational topologies encounters TCP throughput collapse under random non-congestive packet loss over wireless links, high bufferbloat latency spikes, and UDP packet loss in unthrottled UDP video streams.. The complete system architecture must maintain high packet delivery reliability, protect data in transit via cryptographic primitives, and deliver deterministic response latency.

\subsection{System Topology and Engineering Pipeline}
Figure~\ref{fig:pipeline} depicts the modular network architecture and diagnostic pipeline designed for this case study. It spans Ingress Filtering, Deep Packet Inspection, Transport Optimization, Routing Convergence, and Telemetry Monitoring.

\begin{figure}[H]
\centering
\resizebox{\textwidth}{!}{\begin{tikzpicture}[
    block/.style={rectangle, draw=black!80, fill=red!8, thick, rounded corners, minimum height=1.1cm, text width=2.5cm, font=\footnotesize},
    arr/.style={-{Stealth[length=2mm]}, thick}
]
\node[block, align=center] (n1) at (0,0) {1. Client Application\\Bulk File / Live Video};
\node[block, align=center] (n2) at (3.2,0) {2. Transport Layer Mux\\TCP Socket vs UDP Socket};
\node[block, align=center] (n3) at (6.4,0) {3. TCP 3-Way Handshake\\SYN $\to$ SYN-ACK $\to$ ACK};
\node[block, align=center] (n4) at (9.6,0) {4. Congestion Controller\\Reno, Cubic vs BBR};
\node[block, align=center] (n5) at (12.8,0) {5. Sliding Window Flow\\Dynamic Receive Window rcv\_wnd};
\node[block, align=center] (n6) at (12.8,-2.2) {6. UDP Raw Datagram\\Zero-Handshake Fire \& Forget};
\node[block, align=center] (n7) at (9.6,-2.2) {7. WAN Emulation Bridge\\WANem: Delay, Jitter \& Loss};
\node[block, align=center] (n8) at (6.4,-2.2) {8. Wireshark Dissector\\Throughput \& RTT Graphing};
\node[block, align=center] (n9) at (3.2,-2.2) {9. Bufferbloat Tuning\\CoDel / FQ-CoDel Queuing};
\node[block, align=center] (n10) at (0,-2.2) {10. Protocol Selection\\Optimal Transport Mapping};
\draw[arr] (n1) -- (n2);
\draw[arr] (n2) -- (n3);
\draw[arr] (n3) -- (n4);
\draw[arr] (n4) -- (n5);
\draw[arr] (n5) -- (n6);
\draw[arr] (n6) -- (n7);
\draw[arr] (n7) -- (n8);
\draw[arr] (n8) -- (n9);
\draw[arr] (n9) -- (n10);
\end{tikzpicture}}
\caption{System Architecture and Engineering Pipeline for Analysis of TCP and UDP Protocol Performance.}
\label{fig:pipeline}
\end{figure}

\subsection{Architectural Component Specifications}
\begin{enumerate}
  \item \textbf{Border Routing \& ACL Layer:} Filters spoofed IP blocks (BCP 38 anti-spoofing) and applies coarse stateless line-rate access lists at the WAN perimeter.
  \item \textbf{Stateful Inspection \& Threat Engine:} Maintains bi-directional connection tracking tables, inspecting Layer 4 flags and protocol headers for anomalies.
  \item \textbf{Deep Packet Inspection (DPI)::} Snort / Suricata engine parses application-layer payloads against rule databases, extracting malicious shellcode patterns.
  \item \textbf{Core Switching \& Routing Fabric:} Multi-layer Gigabit switches enforce 802.1Q VLAN isolation, Spanning Tree (RSTP), and OSPF area routing.
  \item \textbf{Diagnostic Telemetry Sensor:} Continuous NetFlow/sFlow probes and mirrored SPAN ports deliver real-time metrics to network monitoring systems.
\end{enumerate}
